% ----------------------------------------------------------------------
% Capítulo 2: Referencial teórico
% ----------------------------------------------------------------------
\chapter{Referencial teórico}
\label{cap:referencial}

\instrucao{De oito a quinze páginas. Organize por conceitos, não por autores. Cada seção termina dizendo o que aquele conceito traz para o produto. Cite com \textbackslash cite e \textbackslash citeonline, no sistema autor-data da NBR 10520:2023.}

\section{Desenvolvimento de coleções como processo}
\instrucao{Defina o conceito e descreva as etapas do processo. Uma figura ajuda o leitor a ver o ciclo inteiro.}

O desenvolvimento de coleções é um processo contínuo, em que cada etapa alimenta a seguinte. A Figura~\ref{fig:ciclo} resume as seis etapas mais citadas na literatura da área, da análise da comunidade ao desbaste.

\begin{figure}[htb]
  \centering
  \caption{Etapas do desenvolvimento de coleções}
  \label{fig:ciclo}
  \begin{tikzpicture}[
      etapa/.style={draw=azulbib, fill=azulbib!8, rounded corners, minimum width=3.2cm,
        minimum height=0.9cm, align=center, font=\small},
      seta/.style={-{Stealth}, thick, azulbib}]
    \foreach \nome [count=\i] in {Estudo da\\comunidade, Política de\\seleção, Seleção, Aquisição, Avaliação, Desbaste}{
      \node[etapa] (e\i) at ({90-(\i-1)*60}:3.1cm) {\nome};
    }
    \foreach \i/\j in {1/2,2/3,3/4,4/5,5/6,6/1}{
      \draw[seta] (e\i) -- (e\j);
    }
    \node[font=\small\bfseries, text=azulbib] {Coleção};
  \end{tikzpicture}
  \fonte{Elaborada pelo autor.}
\end{figure}

\section{As leis da Biblioteconomia}
\instrucao{Relacione um princípio clássico com o problema. Exemplo: a segunda lei justifica que a seleção parta do perfil dos leitores.}

As cinco leis formuladas por \citeonline{ranganathan1931} continuam úteis como critério de avaliação de acervos. A segunda lei, \textit{every reader his book}, orienta a seleção a partir do perfil dos leitores. A quinta, \textit{a library is a growing organism}, sustenta a necessidade de prever o crescimento e o desbaste no mesmo documento.

\section{Política de desenvolvimento de coleções}
\instrucao{Liste os elementos mínimos de uma política, de preferência num quadro. Esses elementos viram as seções do produto.}

O Quadro~\ref{qua:elementos} reúne os elementos que o guia deve conter.

\begin{quadro}[htb]
  \caption{Elementos mínimos de uma política de desenvolvimento de coleções}
  \label{qua:elementos}
  \small
  \begin{tabularx}{\textwidth}{lX}
    \toprule
    \textbf{Elemento} & \textbf{O que define} \\
    \midrule
    Objetivos & A missão da biblioteca e o público atendido \\
    Seleção & Critérios de qualidade, pertinência e atualidade \\
    Aquisição & Compra, permuta e doação, com regras para cada via \\
    Avaliação & Indicadores de uso e periodicidade da revisão \\
    Desbaste & Critérios de remanejamento e descarte \\
    \bottomrule
  \end{tabularx}
  \fonte{Elaborado pelo autor.}
\end{quadro}

\section{Produto técnico no mestrado profissional}
\instrucao{Explique o que caracteriza um produto técnico (aplicabilidade, replicabilidade, impacto) e como ele será avaliado. Consulte o regulamento do seu programa.}

No mestrado profissional, a dissertação descreve a pesquisa que fundamenta um produto de uso imediato no ambiente de trabalho. Neste estudo, o produto é um guia com textos-modelo e formulários, que a escola adapta e aprova como política própria.
